\Needspace{14\baselineskip}
\SourceLeaf{102}%
\begin{center}
REMARKS ON THE INVESTIGATIONS INTO THE FUNCTIONS ARISING FROM THE INTEGRAL
$\displaystyle\int\frac{dx}{\BareRoot{(1+x^3)}}$.
\end{center}

\noindent\hspace*{\parindent}According to a diary entry, \textsc{Gauss} inverted the integral considered here
into the concluding power series as early as 9\ September 1796.

The approach developed here bears a relation to \textsc{Weierstrass}’s theory corresponding to
that of the investigation of the lemniscatic integral to \textsc{Jacobi}’s theory.
The latter integral gives the so-called »harmonic« case of the elliptic
structure, while the case here is the »equianharmonic« one.

All of \textsc{Gauss}’s statements follow readily from the addition theorem
for $\boldsymbol{P}(\varphi)$, i.\,e. from the formula given above for
$\boldsymbol{P}(3\psi+\varphi+\varphi')$, which, written somewhat differently, takes the form
\[
\boldsymbol{P}(3\psi+\varphi+\varphi')=
\frac{\begin{gathered}\boldsymbol{P}(\varphi)\boldsymbol{P}(\varphi')
(\boldsymbol{P}(\varphi)+\boldsymbol{P}(\varphi'))\\{}
+2-2\sqrt{1+\boldsymbol{P}(\varphi)^3}\sqrt{1+\boldsymbol{P}(\varphi')^3}\end{gathered}}
{(\boldsymbol{P}(\varphi)-\boldsymbol{P}(\varphi'))^2},
\]
\noindent and which \textsc{Gauss} undoubtedly obtained by the method he also followed for the
lemniscatic integral, i.\,e. through \textsc{Euler}’s original line of reasoning.
Incidentally, this formula reduces directly to the addition theorem for
\textsc{Weierstrass}’s function $\wp(u)$ with $g_2=0$, $g_3=-4$.

\Needspace{12\baselineskip}
The addition theorem readily gives the following formulas.
\[
\boldsymbol{P}(2\varphi+3\psi)=\frac{\boldsymbol{P}(\varphi)(\boldsymbol{P}(\varphi)^3-8)}
{4(\boldsymbol{P}(\varphi)^3+1)},
\]
\[
\boldsymbol{P}(3\varphi)=\frac{\boldsymbol{P}(\varphi)^9-96\boldsymbol{P}(\varphi)^6+48\boldsymbol{P}(\varphi)^3+64}
{9\boldsymbol{P}(\varphi)^2(\boldsymbol{P}(\varphi)^3+4)^2}.
\]
\noindent Using these equations and starting from $\boldsymbol{P}(0)=-1$,
$\boldsymbol{P}(\psi)=0$, one readily derives the remaining statements of \textsc{Gauss}.\hfill\textsc{Fricke.}
